\label{sec:corr:lock}
% Mục 4.2: Kiểm tra mutex lock. Nguồn: tests/test_locks.cpp, results/{apple-m5,C4M7Y670NT,LAP1184}/test.log, tsan.log.

Trong critical section, bài kiểm tra tăng một biến đếm không atomic (đọc, cộng một, ghi) và đếm số
thread đang ở bên trong. Vi phạm hiện ra trực tiếp, hoặc qua một lần cập nhật bị mất. Lượt đầu chạy
số vòng cố định: 1 thread lấy lock 200\,000 lần, 2 thread mỗi thread 100\,000 lần, và số luồng phần cứng
mỗi thread 20\,000 lần. Lượt sau chạy trong 300~ms với 4 thread, số luồng phần cứng và gấp đôi số đó
(\PCname{apple-m5}: 4, 10, 20; hai máy còn lại: 4, 12, 24). Lượt gấp đôi thử lock khi quá tải. Watchdog
kết thúc một bài kiểm tra con sau 20~s và báo deadlock.
% nguồn: tests/test_locks.cpp:220-232 (thread_counts {1, 2, hw}, iter_counts {200000, 100000, 20000},
% timed_threads {4, hw, 2hw}, ms = 300, watchdog 20 s).

Từ số lần lấy lock của từng thread, bài kiểm tra in tỉ số số lần lấy lock lớn nhất / nhỏ nhất (1,00 là
chia đều). Khi mỗi thread có một core, \PCname{mcs} chia gần như đều (1,00 ở 4 thread, 1,01 ở 10 thread),
đúng tính chất FIFO của hàng đợi MCS~\cite{mellorcrummey1991mcs}. Ở 20 thread, MCS chỉ lấy lock 15\,221
lần trong 300~ms (so với 2\,705\,564 lần ở 4 thread), với tỉ số 311,40. Số đo này phù hợp với việc lock
được trao cho một thread đang không có core (cơ chế bàn ở RQ2).
% nguồn: results/apple-m5/test.log, dòng mcs (một lần chạy 300 ms mỗi ô)

Trên \PCname{C4M7Y670NT} (12 core), MCS vẫn chia đều ở 12 thread (tỉ số 1,00, 1\,274\,328 lần lấy lock),
nhưng ở 24 thread chỉ lấy lock 10\,657 lần, với tỉ số 222,25. Trên cả hai máy Apple, MCS chỉ sụp khi số
thread vượt số core.
% nguồn: results/C4M7Y670NT/test.log, dòng mcs (4t 4\,354\,714 acq 1,00; 12t 1\,274\,328 acq 1,00;
% 24t 10\,657 acq 222,25).

Trên \PCname{LAP1184}, MCS sụp sớm hơn. Tỉ số là 1,00 ở 4 thread, nhưng MCS đã sụp ở 12 thread, bằng số
CPU ảo, với 10\,186 lần lấy lock (so với 3\,275\,092 ở 4 thread). Ở 24 thread chỉ còn 1\,628 lần. Ở
12 thread, nhóm cho rằng một CPU ảo đã bị Windows lấy core trong lúc giữ lượt của hàng đợi.
% nguồn: results/LAP1184/test.log, dòng mcs.

Ở bản TSan trên \PCname{LAP1184}, các lượt 100~ms với 24 thread không lấy được lock lần nào (\PCname{mcs}
và \PCname{ttas_bo_spin} ở 12 thread cũng 0). Vì vậy trên máy này TSan chỉ thực sự kiểm các lượt số vòng
cố định và các lượt 4--12 thread.
% nguồn: results/LAP1184/tsan.log, cột 24t = 0 acq ở cả 12 lock; mcs, ttas_bo_spin 12t = 0 acq.

Bản TSan trên \PCname{LAP1184} bỏ qua \PCname{omp_lock} và in ra việc bỏ qua. Khi không bỏ qua, TSan
báo một data race trên biến đếm mà \PCname{omp_lock} bảo vệ. Nhóm cho rằng đây là cảnh báo giả, vì
\texttt{omp\_set\_lock} và \texttt{omp\_unset\_lock} nằm trong libgomp, thư viện không được
instrument. Vì vậy TSan không thấy thứ tự mà lock tạo ra. Cùng lock đó đạt stress test ở bản thường.
% nguồn: results/LAP1184/tsan.log (dòng SKIP omp_lock); lần chạy TSan không bỏ qua (không lưu trong
% results/): SUMMARY: ThreadSanitizer: data race tests/test_locks.cpp:105 in critical_section.
